\subsubsection{Finite formal cone covers and finite exceptional wild inertia}

\begin{lemma}[Proved under stated hypotheses: a finite formal cone cover]
\label{D:finite-formal-cone-cover}
Let $C$ be a smooth proper connected curve in a smooth surface
$Y$ over an algebraically closed field $k$ of characteristic $p>0$.
Write $N=N_{C/Y}$, with $\deg N<0$, and let
\[
 e_Y\in H^1(C,T_C\otimes N^{-1})
\]
be the class of the square-zero algebra extension
$0\to N^{-1}\to\mathcal O_{C_2}\to\mathcal O_C\to0$.
Let $h:C'\to C$ be a finite morphism of smooth proper connected
curves, of degree $m$, and put $N'=h^*N$.  Suppose that
\begin{align}
 h^*e_Y&=0
     &&\text{in }H^1(C',h^*T_C\otimes N'^{-1}),\label{D:cone-first}\\
 H^1(C',h^*T_C\otimes N'^{-n})&=0
     &&(n\geq2),\label{D:cone-tangent}\\
 H^1(C',N'^{\,1-n})&=0
     &&(n\geq2).\label{D:cone-normal}
\end{align}
Then there is a finite locally free formal morphism of rank $m$
\[
 \widehat{\operatorname{Tot}(N')}_{C'}
                 \longrightarrow \widehat Y_C
\]
whose restriction to the reduced closed curves is $h$, and whose
conormal map $h^*N^{-1}\to N'^{-1}$ is the specified identity.
The ramification index along the closed curves is one.

If $Y\to\operatorname{Spec}A$ is a projective resolution of a
normal excellent surface germ, isomorphic off its closed point,
and its reduced exceptional fibre is precisely $C$, then over
$\widehat A$ this formal morphism algebraizes to a finite locally
free morphism $Z\to Y_{\widehat A}$ of rank $m$, with $Z$ regular
and integral.  Its Stein ring $B=H^0(Z,\mathcal O_Z)$ is a finite
normal local $\widehat A$-algebra of fraction-field degree $m$,
and, as a complete local $k$-algebra,
\[
 B\simeq\widehat{\bigoplus_{j\geq0}H^0(C',N'^{-j})}.
 \tag{FC}
\]
The completion is at the positive-degree ideal.  Moreover,
$Z\to\operatorname{Spec}B$ is a resolution, and its exceptional
curve gives the radial valuation of this completed section cone.
No separability or pure inseparability of
$\operatorname{Frac}(B)/\operatorname{Frac}(\widehat A)$ is asserted.
\end{lemma}

\begin{proof}
Put $V=\operatorname{Tot}(N')$, with zero-section ideal $J$,
and $V_n=(C',\mathcal O_V/J^n)$.  Write $I=\mathcal O_Y(-C)$
and $C_n=(C,\mathcal O_Y/I^n)$.  We construct compatible maps
$f_n:V_n\to C_n$ extending $h$ and the stated conormal map.

We first explain the class in \eqref{D:cone-first}.  On affine
opens of $C$, smoothness supplies $k$-algebra splittings of
$\mathcal O_{C_2}\to\mathcal O_C$.  On overlaps their differences
are derivations $\mathcal O_C\to N^{-1}$, giving the cocycle $e_Y$.
Using one such splitting, pullback by $h$ and the prescribed
map of the square-zero ideals give a local algebra map
$\mathcal O_{C_2}\to h_*\mathcal O_{V_2}$.  The differences of
these maps are the pulled-back cocycle, now in
$\mathcal Hom(h^*\Omega_C,N'^{-1})$ on $C'$.
Changing a local map by a derivation changes this cocycle by
a coboundary.  Thus \eqref{D:cone-first} gives a global $f_2$.
This argument does not assume a global retraction $C_2\to C$,
and it does not identify the prescribed conormal map with the
zero conormal map of absolute Frobenius.

For $n\geq2$, regard $f_n$ as a morphism $V_n\to Y$.
The square-zero ideal of $V_n$ in $V_{n+1}$ is
$J^n/J^{n+1}=N'^{-n}$.  Smoothness of $Y$ gives local lifts,
and their differences form the sheaf
\[
 \mathcal H_n=h^*(T_Y|_C)\otimes N'^{-n}.
 \tag{FL}
\]
The obstruction is its first cohomology.  The tangent-normal
sequence, pulled back along the flat map $h$, is exact and gives
\[
 0\longrightarrow h^*T_C\otimes N'^{-n}
 \longrightarrow\mathcal H_n
 \longrightarrow N'^{\,1-n}\longrightarrow0.
 \tag{FN}
\]
A finite nonconstant map between smooth curves is flat, as
its local modules are torsion free over DVRs.  Conditions
\eqref{D:cone-tangent}--\eqref{D:cone-normal} therefore imply
$H^1(C',\mathcal H_n)=0$ and give a global lift $f_{n+1}$.
It factors through $C_{n+1}$ because it sends $I$ into $J$.
Its conormal map remains that of $f_2$.
The precise square-zero lifting and derivation inputs are
\cite[Lemmas 17.28.11, 37.9.4--5 and 37.11.7,
tags 01UP, 04FH, 04FJ and 02H6]{StacksBNegativeCurve}.

Taking inverse limits gives the claimed formal map.  We check
finiteness and flatness, rather than inferring them solely from
the finite closed-fibre map.  Locally on the target curve choose
a generator $t$ of $I$ and choose $s$ for $J$ on its inverse
image.  The image of $t$ is $s$ times a unit, by the prescribed
conormal isomorphism.  Modulo this parameter the map is $h$,
which is finite locally free of rank $m$.  After restricting
the target open, lift a basis of that finite free module.
Successive reduction modulo $t^r$ and completeness show that
the lifted basis generates the formal algebra.  A relation
between its members reduces to zero modulo $t$, hence all its
coefficients are divisible by $t$.  Both formal rings are
$t$-torsion free; repeating and using separatedness shows that
the relation was zero.  Thus the formal algebra is locally free
of rank $m$.  This also proves the stated ramification index.

For the algebraization assertion, replace $A$ by its completion.
The $\mathfrak m_A$-adic and $I$-adic topologies on the proper
resolution are equivalent: their defining ideals have the same
radical, and a finite affine covering gives cofinal powers.
Our finite formal algebra consequently gives a compatible
system of finite algebras on $Y\times_A A/\mathfrak m_A^r$.
The precise finite-morphism form of Grothendieck existence,
\cite[Lemma 30.28.2, tag 09ZT]{StacksBFiniteCone},
algebraizes it to $Z\to Y$.  In that theorem the closed support
is proper, and its proof algebraizes the multiplication and
unit using full faithfulness; no separate algebraization of
a nonproper formal neighborhood is required.

Flatness and rank $m$ hold along the closed fibre, by faithful
flatness of completed local rings.  The nonflat locus and the
loci of other ranks are closed.  If nonempty they would have
a closed point over the closed point of $\operatorname{Spec}A$,
since $Z$ and $Y$ are proper over this local base.  They are
therefore empty.  At every point of the closed fibre the
completed local ring of $Z$ agrees with that of the smooth
line-bundle model $V$; hence $Z$ is regular there.  Excellence
and the same properness argument make $Z$ regular everywhere.
Its closed fibre is connected, with reduction $C'$; every
connected component of $Z$ meets it.  Thus $Z$ is connected,
and a connected regular scheme is integral in this situation.

Proper coherent cohomology makes $B=H^0(Z,\mathcal O_Z)$ finite
over $A$.  Formal functions
\cite[Lemma 30.20.6, tag 087U]{StacksBFiniteCone}
identifies it with
\[
 H^0(\widehat V_{C'},\mathcal O)
  =\varprojlim_r\bigoplus_{0\leq j<r}H^0(C',N'^{-j})
  =\prod_{j\geq0}H^0(C',N'^{-j}).
 \tag{FS}
\]
The section ring of the ample line bundle $N'^{-1}$ is a
finitely generated normal positive graded domain with degree
zero part $k$.  As in Proposition~\ref{D:graded-cone-tame},
finite generation makes its positive-ideal topology equivalent
to the degree topology.  Thus (FS) is exactly (FC).
In particular $B$ is local and normal.  Normality also follows
directly because a global function on an integral normal
scheme that is integral over its global ring lies in every
local ring and hence is global.

The positive-rank finite locally free morphism $Z\to Y$ is
surjective, so the map $A\to B$ is injective.
On the puncture of $\operatorname{Spec}A$ the original resolution
is an isomorphism, so $Z$ is finite there.  Its Stein morphism
to $\operatorname{Spec}B$ is consequently an isomorphism off
the closed point.  On the generic fibre, finite flatness gives
a field of dimension $m$ over $\operatorname{Frac}A$; flat base
change of $H^0$ identifies it with
$B\otimes_A\operatorname{Frac}A$.  This proves the degree claim
and birationality of the Stein morphism.  At the generic point
of $C'$, a series in (FS) has valuation its least nonzero
homogeneous degree.  This is precisely the radial valuation
of the section cone, and it restricts to the valuation of $C$
with index one.
\end{proof}

\begin{proposition}[Proved under stated hypotheses: a fixed cone cover bounds all wild images]
\label{D:finite-cone-wild-bound}
In the algebraic situation of Lemma~\ref{D:finite-formal-cone-cover},
the image of the wild inertia of $C$ in the punctured
fundamental group of the surface germ is finite of order at
most $m$.  The same bound holds after strict henselization.
It bounds the wild image in every finite quasi-\'etale cover,
with no solvability or degree restriction on the covering group.
\end{proposition}

\begin{proof}
First work over the complete local ring $A=\widehat A$ and
use the finite local normal cone algebra $B$ constructed above.
The inverse image of the closed point is its unique closed
point.  Thus every finite \'etale cover of the puncture of $A$
pulls back to a finite \'etale cover of the puncture of $B$.
Proposition~\ref{D:graded-cone-tame} makes every such pulled-back
cover tame at the radial valuation $w$ of $B$.

Let $v$ be the valuation of $C$.  The actual finite morphism
$Z\to Y_A$ identifies $w$ as an extension of $v$; the corresponding
complete DVR fraction fields form a finite extension $L/K$ of
degree at most $m$.  We record the wild-inertia comparison
including inseparable extensions.  For a finite separable
extension, compatible choices inside a separable closure give
\[
 I_L=I_K\cap G_L,\qquad P_L=P_K\cap G_L. \tag{WI}
\]
Indeed the first equality follows from the action on the
common algebraic residue closure.  Inertia has a normal
pro-$p$ subgroup with pro-prime-to-$p$ quotient, by
\cite[Lemma 58.13.5, tag 0BUA]{StacksBFiniteCone}.
Intersecting with $G_L$ makes $P_K\cap G_L$ the pro-$p$
Sylow subgroup of $I_L$, which is its wild inertia.
For a finite purely inseparable extension, the equivalence
of separable closures and the uniquely extended valuation
identify the absolute Galois groups, their inertia subgroups,
and their characteristic wild subgroups.  Factoring a general
finite extension into its maximal separable part and a purely
inseparable extension proves the same assertion with
$G_L$ interpreted under this identification.

The preceding tame pullback kills $P_L$ in the monodromy of
every original punctured cover.  Thus, for any such monodromy
representation $\rho$,
\[
 |\rho(P_K)|
 \leq [P_K:P_K\cap G_L]
 \leq [L:K]_{\rm sep}
 \leq m. \tag{WB}
\]
The subgroup $G_L$ need not be normal.  All representations
kill the same intersection in (WI), so the full image of
$P_K$ in the punctured fundamental group is already a
quotient of a set of at most $m$ cosets; equivalently its
finite quotients have uniformly bounded order.  It is a
finite pro-$p$ group of order at most $m$.

Finally suppose the original germ is not complete.  Base
change the proper resolution to the completed local ring.
At the generic point of its closed exceptional curve, the
original DVR and the new DVR have isomorphic genuine DVR
completions, by \cite[Lemma 54.11.1, tag 0BG4]
{StacksBFiniteCone}.  Normalize in any finite cover and make
the same comparison.  The resulting finite extensions of
complete DVRs, and hence their wild inertia images, are the
ones just bounded.  This argument applies to each component
if a cover splits after completion.  Strict henselization
has the same completion here, since the residue field $k$ is
algebraically closed; alternatively each of its finite covers
descends to an \'etale neighborhood with that same completed
resolution.  Thus (WB) applies in all the stated versions.
We have not identified completion at the surface singularity
with completion at its boundary valuation.
\end{proof}

\begin{corollary}[Proved under stated hypotheses: the threshold $d>2g-2$]
\label{D:negative-curve-finite-wild}
Let a normal surface germ over an algebraically closed field
of characteristic $p>0$ admit a projective resolution with
sole exceptional curve a smooth proper connected curve $C$,
of genus $g$ and self-intersection $-d<0$.
If
\[
                         d>\max\{0,2g-2\},
\]
the exceptional wild inertia has finite image in its full
punctured fundamental group.  More precisely choose a power
$q=p^r$ satisfying
\[
 q(d-2g+2)>2g-2,\qquad qd>2g-2.
 \tag{Q}
\]
Its image has order at most $q$.
\end{corollary}

\begin{proof}
Take the $k$-linear relative Frobenius
$h:C^{(-r)}\to C$, of degree $q$, and put $C'=C^{(-r)}$.
The genus remains $g$.  The line bundle in the first
obstruction has degree $q(d-2g+2)>2g-2$, so its $H^1$
vanishes by Serre duality.  The two line bundles in the
later obstruction have degrees
\[
 q(nd-2g+2),\qquad q(n-1)d,\qquad n\geq2.
 \tag{FD}
\]
Both exceed $2g-2$ by (Q).  Apply the lemma and proposition.
The degree and Serre-duality inputs are
\cite[Lemmas 53.4.1 and 53.5.2, tags 0BS2 and 0BS6]
{StacksBNegativeCurve}.  The conclusion is finiteness of
the full exceptional wild image; it need not be trivial.
\end{proof}
